
This paper measures semantic accommodation in human-AI conversations drawn from the Anthropic HH-RLHF helpfulness dataset, where conversations are stratified across three reinforcement learning from human feedback (RLHF) alignment quality tiers. We introduce C\_sem, a training-free sentence-BERT (SBERT) metric that quantifies interaction-specific semantic accommodation by subtracting a random partner-shuffle baseline from turn-pair cosine similarity. A three-level partner-specificity hierarchy (actual partner > KNN topic-matched > random) provides additional validation that the measured signal exceeds topically matched controls. Analysis of 155,362 conversation pairs indicates that human follow-up turns exhibit substantial semantic similarity to their actual AI partner turns relative to random baseline turns (C\_sem = 0.329, 95\% CI [0.328, 0.330]; Cohen's d = 1.998 vs. random; d = 0.417 vs. KNN topic-matched). C\_sem increases monotonically across RLHF quality tiers (Jonckheere-Terpstra p = 0.001, replicated across three SBERT architectures; Cohen's d for T1 to T3 ranges from 0.183 to 0.254). In all nine tier-by-model conditions, human-to-AI accommodation exceeds AI-to-human accommodation (d = 0.061--0.41). However, two mechanism hypotheses are falsified: within-conversation quality discrimination yields a reversed signal (human follow-up turns are more similar to rejected AI responses than to chosen ones in 25 of 27 conditions), and a politeness-style mediation proxy shows no predictive power (beta approximately 0, p approximately 0.99). These results characterize the accommodation pattern as a population-level association between RLHF tier and human semantic behavior, rather than a within-conversation perceptual effect. The cross-sectional design precludes causal claims.
